%% Generated by Sphinx.
\def\sphinxdocclass{report}
\documentclass[letterpaper,10pt,english]{sphinxmanual}
\ifdefined\pdfpxdimen
   \let\sphinxpxdimen\pdfpxdimen\else\newdimen\sphinxpxdimen
\fi \sphinxpxdimen=.75bp\relax
\ifdefined\pdfimageresolution
    \pdfimageresolution= \numexpr \dimexpr1in\relax/\sphinxpxdimen\relax
\fi
\newdimen\sphinxremdimen\sphinxremdimen = 10pt
%% let collapsible pdf bookmarks panel have high depth per default
\PassOptionsToPackage{bookmarksdepth=5}{hyperref}

\PassOptionsToPackage{booktabs}{sphinx}
\PassOptionsToPackage{colorrows}{sphinx}

\PassOptionsToPackage{warn}{textcomp}
\usepackage[utf8]{inputenc}
\ifdefined\DeclareUnicodeCharacter
% support both utf8 and utf8x syntaxes
  \ifdefined\DeclareUnicodeCharacterAsOptional
    \def\sphinxDUC#1{\DeclareUnicodeCharacter{"#1}}
  \else
    \let\sphinxDUC\DeclareUnicodeCharacter
  \fi
  \sphinxDUC{00A0}{\nobreakspace}
  \sphinxDUC{2500}{\sphinxunichar{2500}}
  \sphinxDUC{2502}{\sphinxunichar{2502}}
  \sphinxDUC{2514}{\sphinxunichar{2514}}
  \sphinxDUC{251C}{\sphinxunichar{251C}}
  \sphinxDUC{2572}{\textbackslash}
\fi
\usepackage{cmap}
\usepackage[T1]{fontenc}
\usepackage{amsmath,amssymb,amstext}
\usepackage{babel}



\usepackage{tgtermes}
\usepackage{tgheros}
\renewcommand{\ttdefault}{txtt}



\usepackage[Bjarne]{fncychap}
\usepackage{sphinx}

\fvset{fontsize=auto}
\usepackage{geometry}


% Include hyperref last.
\usepackage{hyperref}
% Fix anchor placement for figures with captions.
\usepackage{hypcap}% it must be loaded after hyperref.
% Set up styles of URL: it should be placed after hyperref.
\urlstyle{same}

\addto\captionsenglish{\renewcommand{\contentsname}{Project Documentation}}

\usepackage{sphinxmessages}
\setcounter{tocdepth}{1}



\title{User Access Control}
\date{Nov 27, 2026}
\release{1.0}
\author{Project Author}
\newcommand{\sphinxlogo}{\vbox{}}
\renewcommand{\releasename}{Release}
\makeindex
\begin{document}

\ifdefined\shorthandoff
  \ifnum\catcode`\=\string=\active\shorthandoff{=}\fi
  \ifnum\catcode`\"=\active\shorthandoff{"}\fi
\fi

\pagestyle{empty}
\sphinxmaketitle
\pagestyle{plain}
\sphinxtableofcontents
\pagestyle{normal}
\phantomsection\label{\detokenize{index::doc}}


\sphinxAtStartPar
Welcome to the documentation of the User Access Control project.

\sphinxAtStartPar
This project is developed using Python and FastAPI. It provides a secure
authentication and authorization system that allows users to log in and
access only the pages permitted to them.


\chapter{Project Features}
\label{\detokenize{index:project-features}}\begin{itemize}
\item {} 
\sphinxAtStartPar
User login and logout

\item {} 
\sphinxAtStartPar
Password hashing using bcrypt

\item {} 
\sphinxAtStartPar
Session\sphinxhyphen{}based authentication

\item {} 
\sphinxAtStartPar
User\sphinxhyphen{}specific page permissions

\item {} 
\sphinxAtStartPar
Protected routes using FastAPI

\item {} 
\sphinxAtStartPar
MySQL database integration

\item {} 
\sphinxAtStartPar
Access denied page for unauthorized users

\end{itemize}


\chapter{Technology Stack}
\label{\detokenize{index:technology-stack}}\begin{itemize}
\item {} 
\sphinxAtStartPar
Python

\item {} 
\sphinxAtStartPar
FastAPI

\item {} 
\sphinxAtStartPar
MySQL

\item {} 
\sphinxAtStartPar
PyMySQL

\item {} 
\sphinxAtStartPar
Starlette SessionMiddleware

\item {} 
\sphinxAtStartPar
bcrypt

\item {} 
\sphinxAtStartPar
HTML and CSS

\end{itemize}


\chapter{Documentation Contents}
\label{\detokenize{index:documentation-contents}}
\sphinxstepscope


\section{Project Overview}
\label{\detokenize{overview:project-overview}}\label{\detokenize{overview::doc}}

\subsection{Introduction}
\label{\detokenize{overview:introduction}}
\sphinxAtStartPar
User Access Control is a web application developed using Python and FastAPI.
It provides authentication and authorization functionality to control which
pages a logged\sphinxhyphen{}in user can access.

\sphinxAtStartPar
The application verifies user credentials during login and uses session\sphinxhyphen{}based
authentication to maintain the user’s login state.


\subsection{Objectives}
\label{\detokenize{overview:objectives}}\begin{itemize}
\item {} 
\sphinxAtStartPar
Provide secure user authentication.

\item {} 
\sphinxAtStartPar
Store user information in a MySQL database.

\item {} 
\sphinxAtStartPar
Protect application pages from unauthorized access.

\item {} 
\sphinxAtStartPar
Assign different page permissions to different users.

\item {} 
\sphinxAtStartPar
Display an access denied page when a user tries to access a restricted page.

\end{itemize}


\subsection{Main Features}
\label{\detokenize{overview:main-features}}

\subsubsection{User Authentication}
\label{\detokenize{overview:user-authentication}}
\sphinxAtStartPar
Users can log in using their username and password. Passwords are stored
as bcrypt hashes rather than plain text.


\subsubsection{Session Management}
\label{\detokenize{overview:session-management}}
\sphinxAtStartPar
The application uses Starlette SessionMiddleware to maintain user login
sessions.


\subsubsection{Page\sphinxhyphen{}Level Authorization}
\label{\detokenize{overview:page-level-authorization}}
\sphinxAtStartPar
Each user is assigned permission to access specific pages. The application
checks these permissions before displaying protected pages.


\subsection{Technology Used}
\label{\detokenize{overview:technology-used}}\begin{itemize}
\item {} 
\sphinxAtStartPar
Python 3.13

\item {} 
\sphinxAtStartPar
FastAPI

\item {} 
\sphinxAtStartPar
MySQL

\item {} 
\sphinxAtStartPar
PyMySQL

\item {} 
\sphinxAtStartPar
Starlette

\item {} 
\sphinxAtStartPar
bcrypt

\item {} 
\sphinxAtStartPar
HTML

\item {} 
\sphinxAtStartPar
CSS

\end{itemize}


\subsection{Application Pages}
\label{\detokenize{overview:application-pages}}
\sphinxAtStartPar
The application includes the following protected pages:
\begin{itemize}
\item {} 
\sphinxAtStartPar
Page 1

\item {} 
\sphinxAtStartPar
Page 2

\item {} 
\sphinxAtStartPar
Page 3

\item {} 
\sphinxAtStartPar
Page 4

\end{itemize}

\sphinxAtStartPar
The pages available to a user depend on the permissions assigned to that user.

\sphinxstepscope


\section{Installation Guide}
\label{\detokenize{installation:installation-guide}}\label{\detokenize{installation::doc}}

\subsection{Prerequisites}
\label{\detokenize{installation:prerequisites}}
\sphinxAtStartPar
Before running the application, make sure the following software is installed:
\begin{itemize}
\item {} 
\sphinxAtStartPar
Python 3.13

\item {} 
\sphinxAtStartPar
MySQL Server

\item {} 
\sphinxAtStartPar
MySQL client or MySQL Workbench

\item {} 
\sphinxAtStartPar
Visual Studio Code or another code editor

\end{itemize}


\subsection{Download the Project}
\label{\detokenize{installation:download-the-project}}
\sphinxAtStartPar
Open the project folder in your terminal.


\subsection{Install Required Packages}
\label{\detokenize{installation:install-required-packages}}
\sphinxAtStartPar
Create a virtual environment:

\begin{sphinxVerbatim}[commandchars=\\\{\}]
python\PYG{+w}{ }\PYGZhy{}m\PYG{+w}{ }venv\PYG{+w}{ }venv
\end{sphinxVerbatim}

\sphinxAtStartPar
Activate the virtual environment on Windows:

\begin{sphinxVerbatim}[commandchars=\\\{\}]
venv\PYGZbs{}Scripts\PYGZbs{}activate
\end{sphinxVerbatim}

\sphinxAtStartPar
Install the required Python packages:

\begin{sphinxVerbatim}[commandchars=\\\{\}]
python\PYG{+w}{ }\PYGZhy{}m\PYG{+w}{ }pip\PYG{+w}{ }install\PYG{+w}{ }fastapi\PYG{+w}{ }uvicorn\PYG{+w}{ }pymysql\PYG{+w}{ }bcrypt\PYG{+w}{ }python\PYGZhy{}dotenv\PYG{+w}{ }itsdangerous
\end{sphinxVerbatim}


\subsection{Database Configuration}
\label{\detokenize{installation:database-configuration}}
\sphinxAtStartPar
Make sure MySQL Server is running.

\sphinxAtStartPar
Create the database and configure the database connection according to the
settings used by the application’s database module.


\subsection{Environment Configuration}
\label{\detokenize{installation:environment-configuration}}
\sphinxAtStartPar
Configure the required environment variables, such as database credentials
and session secret, according to the project’s configuration.


\subsection{Run the Application}
\label{\detokenize{installation:run-the-application}}
\sphinxAtStartPar
Start the FastAPI application using:

\begin{sphinxVerbatim}[commandchars=\\\{\}]
python\PYG{+w}{ }\PYGZhy{}m\PYG{+w}{ }uvicorn\PYG{+w}{ }app.main:app\PYG{+w}{ }\PYGZhy{}\PYGZhy{}reload
\end{sphinxVerbatim}

\sphinxAtStartPar
Open the following URL in your browser:

\begin{sphinxVerbatim}[commandchars=\\\{\}]
http://127.0.0.1:8000
\end{sphinxVerbatim}


\subsection{API Documentation}
\label{\detokenize{installation:api-documentation}}
\sphinxAtStartPar
FastAPI provides interactive API documentation at:

\begin{sphinxVerbatim}[commandchars=\\\{\}]
http://127.0.0.1:8000/docs
\end{sphinxVerbatim}

\sphinxstepscope


\section{System Architecture}
\label{\detokenize{architecture:system-architecture}}\label{\detokenize{architecture::doc}}

\subsection{Introduction}
\label{\detokenize{architecture:introduction}}
\sphinxAtStartPar
The User Access Control application follows a modular architecture.
The application separates authentication, authorization, database operations,
security, and utility functions into different Python modules.


\subsection{Application Structure}
\label{\detokenize{architecture:application-structure}}
\sphinxAtStartPar
The project contains the following main components:

\begin{sphinxVerbatim}[commandchars=\\\{\}]
UserAccessControl/
|
+\PYGZhy{}\PYGZhy{} app/
|   |
|   +\PYGZhy{}\PYGZhy{} main.py
|   +\PYGZhy{}\PYGZhy{} auth.py
|   +\PYGZhy{}\PYGZhy{} access.py
|   +\PYGZhy{}\PYGZhy{} db.py
|   +\PYGZhy{}\PYGZhy{} security.py
|   +\PYGZhy{}\PYGZhy{} utils.py
|   |
|   +\PYGZhy{}\PYGZhy{} html/
|   |   +\PYGZhy{}\PYGZhy{} login.html
|   |   +\PYGZhy{}\PYGZhy{} home.html
|   |   +\PYGZhy{}\PYGZhy{} denied.html
|   |   +\PYGZhy{}\PYGZhy{} page1.html
|   |   +\PYGZhy{}\PYGZhy{} page2.html
|   |   +\PYGZhy{}\PYGZhy{} page3.html
|   |   +\PYGZhy{}\PYGZhy{} page4.html
|   |
|   +\PYGZhy{}\PYGZhy{} static/
|
+\PYGZhy{}\PYGZhy{} docs/
    +\PYGZhy{}\PYGZhy{} source/
    +\PYGZhy{}\PYGZhy{} build/
\end{sphinxVerbatim}


\subsection{Main Components}
\label{\detokenize{architecture:main-components}}

\subsubsection{main.py}
\label{\detokenize{architecture:main-py}}
\sphinxAtStartPar
The main application module creates the FastAPI application, configures
session middleware, mounts static files, initializes the database, and
registers application routes.


\subsubsection{auth.py}
\label{\detokenize{architecture:auth-py}}
\sphinxAtStartPar
This module handles user authentication, including login and logout
functionality.


\subsubsection{access.py}
\label{\detokenize{architecture:access-py}}
\sphinxAtStartPar
This module manages user\sphinxhyphen{}specific page permissions and controls access
to protected pages.


\subsubsection{db.py}
\label{\detokenize{architecture:db-py}}
\sphinxAtStartPar
This module handles database connectivity and database\sphinxhyphen{}related operations.


\subsubsection{security.py}
\label{\detokenize{architecture:security-py}}
\sphinxAtStartPar
This module contains security\sphinxhyphen{}related functionality, such as password
hashing and password verification.


\subsubsection{utils.py}
\label{\detokenize{architecture:utils-py}}
\sphinxAtStartPar
This module contains reusable helper functions used by the application.


\subsection{Architecture Diagram}
\label{\detokenize{architecture:architecture-diagram}}
\sphinxAtStartPar
The diagram below shows how a browser request moves through the
application components.


\subsection{Module Dependency Diagram}
\label{\detokenize{architecture:module-dependency-diagram}}
\sphinxAtStartPar
The diagram below shows which modules depend on which other modules.


\subsection{Authentication and Authorization Flow}
\label{\detokenize{architecture:authentication-and-authorization-flow}}
\sphinxAtStartPar
The application follows this general workflow:
\begin{enumerate}
\sphinxsetlistlabels{\arabic}{enumi}{enumii}{}{.}%
\item {} 
\sphinxAtStartPar
The user opens the login page.

\item {} 
\sphinxAtStartPar
The user submits their username and password.

\item {} 
\sphinxAtStartPar
The application verifies the credentials.

\item {} 
\sphinxAtStartPar
After successful authentication, a user session is created.

\item {} 
\sphinxAtStartPar
The user opens a protected page.

\item {} 
\sphinxAtStartPar
The application checks the user’s page permissions.

\item {} 
\sphinxAtStartPar
If permission is granted, the requested page is displayed.

\item {} 
\sphinxAtStartPar
Otherwise, the access denied page is displayed.

\end{enumerate}


\subsubsection{Login Sequence Diagram}
\label{\detokenize{architecture:login-sequence-diagram}}

\subsubsection{Page Access Sequence Diagram}
\label{\detokenize{architecture:page-access-sequence-diagram}}

\subsection{Conclusion}
\label{\detokenize{architecture:conclusion}}
\sphinxAtStartPar
The modular architecture makes the application easier to understand,
maintain, test, and extend.

\sphinxstepscope


\section{Database Documentation}
\label{\detokenize{database:database-documentation}}\label{\detokenize{database::doc}}

\subsection{Introduction}
\label{\detokenize{database:introduction}}
\sphinxAtStartPar
The User Access Control application uses MySQL as its relational database
management system.

\sphinxAtStartPar
The database stores user information and page permission details required
for authentication and authorization.


\subsection{Database Connectivity}
\label{\detokenize{database:database-connectivity}}
\sphinxAtStartPar
The application uses PyMySQL to connect Python code with the MySQL database.

\sphinxAtStartPar
The database module is responsible for managing database connections and
executing database\sphinxhyphen{}related operations.


\subsection{Database Tables}
\label{\detokenize{database:database-tables}}

\subsubsection{Users Table}
\label{\detokenize{database:users-table}}
\sphinxAtStartPar
The users table stores registered user information.

\sphinxAtStartPar
Typical information includes:
\begin{itemize}
\item {} 
\sphinxAtStartPar
User ID

\item {} 
\sphinxAtStartPar
Username

\item {} 
\sphinxAtStartPar
Full name

\item {} 
\sphinxAtStartPar
Password hash

\end{itemize}


\subsubsection{User Permissions Table}
\label{\detokenize{database:user-permissions-table}}
\sphinxAtStartPar
The user\_pages table stores the pages that each user is allowed to access.

\sphinxAtStartPar
Typical information includes:
\begin{itemize}
\item {} 
\sphinxAtStartPar
User ID

\item {} 
\sphinxAtStartPar
Page name

\end{itemize}


\subsection{User Permission Relationship}
\label{\detokenize{database:user-permission-relationship}}
\sphinxAtStartPar
A user can have permission to access multiple pages.

\sphinxAtStartPar
For example:

\begin{sphinxVerbatim}[commandchars=\\\{\}]
Username    Allowed Pages
\PYGZhy{}\PYGZhy{}\PYGZhy{}\PYGZhy{}\PYGZhy{}\PYGZhy{}\PYGZhy{}\PYGZhy{}    \PYGZhy{}\PYGZhy{}\PYGZhy{}\PYGZhy{}\PYGZhy{}\PYGZhy{}\PYGZhy{}\PYGZhy{}\PYGZhy{}\PYGZhy{}\PYGZhy{}\PYGZhy{}\PYGZhy{}\PYGZhy{}\PYGZhy{}\PYGZhy{}\PYGZhy{}\PYGZhy{}\PYGZhy{}\PYGZhy{}
priti       page1, page2
yogesh      page3
amit        page4
admin       page1, page2, page3, page4
\end{sphinxVerbatim}


\subsection{Database Operations}
\label{\detokenize{database:database-operations}}
\sphinxAtStartPar
The application performs the following database operations:
\begin{itemize}
\item {} 
\sphinxAtStartPar
Establishing database connections.

\item {} 
\sphinxAtStartPar
Retrieving user information.

\item {} 
\sphinxAtStartPar
Verifying user credentials.

\item {} 
\sphinxAtStartPar
Retrieving page permissions.

\item {} 
\sphinxAtStartPar
Storing user and permission information.

\end{itemize}


\subsection{Database Security}
\label{\detokenize{database:database-security}}
\sphinxAtStartPar
The application should follow these database security practices:
\begin{itemize}
\item {} 
\sphinxAtStartPar
Use parameterized SQL queries.

\item {} 
\sphinxAtStartPar
Do not store plain\sphinxhyphen{}text passwords.

\item {} 
\sphinxAtStartPar
Keep database credentials in environment variables.

\item {} 
\sphinxAtStartPar
Restrict database access to authorized users.

\item {} 
\sphinxAtStartPar
Close database connections after use.

\end{itemize}


\subsection{Conclusion}
\label{\detokenize{database:conclusion}}
\sphinxAtStartPar
The MySQL database provides persistent storage for user information and
page\sphinxhyphen{}level permissions.

\sphinxstepscope


\section{Authentication}
\label{\detokenize{authentication:authentication}}\label{\detokenize{authentication::doc}}

\subsection{Introduction}
\label{\detokenize{authentication:introduction}}
\sphinxAtStartPar
Authentication is the process of verifying the identity of a user before
allowing access to the application.

\sphinxAtStartPar
The User Access Control application uses username and password authentication
with bcrypt password hashing.


\subsection{Login Process}
\label{\detokenize{authentication:login-process}}
\sphinxAtStartPar
The login process follows these steps:
\begin{enumerate}
\sphinxsetlistlabels{\arabic}{enumi}{enumii}{}{.}%
\item {} 
\sphinxAtStartPar
The user opens the login page.

\item {} 
\sphinxAtStartPar
The user enters a username and password.

\item {} 
\sphinxAtStartPar
The application receives the submitted credentials.

\item {} 
\sphinxAtStartPar
The application retrieves the corresponding user information from MySQL.

\item {} 
\sphinxAtStartPar
The submitted password is checked against the stored bcrypt password hash.

\item {} 
\sphinxAtStartPar
If the credentials are valid, the application creates a user session.

\item {} 
\sphinxAtStartPar
The user is redirected to the home page.

\item {} 
\sphinxAtStartPar
If the credentials are invalid, an error message is displayed.

\end{enumerate}


\subsection{Login Sequence Diagram}
\label{\detokenize{authentication:login-sequence-diagram}}

\subsection{Password Hashing}
\label{\detokenize{authentication:password-hashing}}
\sphinxAtStartPar
The application uses bcrypt to protect user passwords.

\sphinxAtStartPar
Password hashing converts a password into a secure hash before storing it
in the database.

\sphinxAtStartPar
During login, bcrypt verifies the submitted password against the stored hash.


\subsection{Code Example}
\label{\detokenize{authentication:code-example}}
\sphinxAtStartPar
The following example demonstrates password hashing and password verification.

\begin{sphinxVerbatim}[commandchars=\\\{\}]
\PYG{k+kn}{from}\PYG{+w}{ }\PYG{n+nn}{app}\PYG{n+nn}{.}\PYG{n+nn}{security}\PYG{+w}{ }\PYG{k+kn}{import} \PYG{n}{hash\PYGZus{}password}\PYG{p}{,} \PYG{n}{verify\PYGZus{}password}

\PYG{c+c1}{\PYGZsh{} Generate a password hash during user registration}
\PYG{n}{stored\PYGZus{}hash} \PYG{o}{=} \PYG{n}{hash\PYGZus{}password}\PYG{p}{(}\PYG{l+s+s2}{\PYGZdq{}}\PYG{l+s+s2}{MySecret@123}\PYG{l+s+s2}{\PYGZdq{}}\PYG{p}{)}

\PYG{c+c1}{\PYGZsh{} Verify the password during login}
\PYG{n}{is\PYGZus{}valid} \PYG{o}{=} \PYG{n}{verify\PYGZus{}password}\PYG{p}{(}\PYG{l+s+s2}{\PYGZdq{}}\PYG{l+s+s2}{MySecret@123}\PYG{l+s+s2}{\PYGZdq{}}\PYG{p}{,} \PYG{n}{stored\PYGZus{}hash}\PYG{p}{)}

\PYG{k}{if} \PYG{n}{is\PYGZus{}valid}\PYG{p}{:}
    \PYG{n+nb}{print}\PYG{p}{(}\PYG{l+s+s2}{\PYGZdq{}}\PYG{l+s+s2}{Login successful}\PYG{l+s+s2}{\PYGZdq{}}\PYG{p}{)}
\PYG{k}{else}\PYG{p}{:}
    \PYG{n+nb}{print}\PYG{p}{(}\PYG{l+s+s2}{\PYGZdq{}}\PYG{l+s+s2}{Invalid password}\PYG{l+s+s2}{\PYGZdq{}}\PYG{p}{)}
\end{sphinxVerbatim}


\subsection{Session Management}
\label{\detokenize{authentication:session-management}}
\sphinxAtStartPar
The application uses Starlette SessionMiddleware to manage login sessions.

\sphinxAtStartPar
After successful login, the user’s ID is stored in the session.

\sphinxAtStartPar
The session middleware manages the session cookie in the user’s browser.


\subsubsection{Session Configuration}
\label{\detokenize{authentication:session-configuration}}
\sphinxAtStartPar
The application configures session settings such as:
\begin{itemize}
\item {} 
\sphinxAtStartPar
Session secret key

\item {} 
\sphinxAtStartPar
SameSite cookie policy

\item {} 
\sphinxAtStartPar
HTTPS\sphinxhyphen{}only cookie setting

\item {} 
\sphinxAtStartPar
Session expiration time

\end{itemize}


\subsection{Logout Process}
\label{\detokenize{authentication:logout-process}}
\sphinxAtStartPar
The logout process follows these steps:
\begin{enumerate}
\sphinxsetlistlabels{\arabic}{enumi}{enumii}{}{.}%
\item {} 
\sphinxAtStartPar
The user selects the logout option.

\item {} 
\sphinxAtStartPar
The application clears the user’s session.

\item {} 
\sphinxAtStartPar
The user is redirected to the login page.

\end{enumerate}


\subsection{Authentication Security}
\label{\detokenize{authentication:authentication-security}}
\sphinxAtStartPar
The following security practices are important:
\begin{itemize}
\item {} 
\sphinxAtStartPar
Store password hashes instead of plain\sphinxhyphen{}text passwords.

\item {} 
\sphinxAtStartPar
Use a strong session secret.

\item {} 
\sphinxAtStartPar
Configure secure cookies when deploying with HTTPS.

\item {} 
\sphinxAtStartPar
Validate user input.

\item {} 
\sphinxAtStartPar
Protect login routes against repeated password attempts.

\end{itemize}


\subsection{Conclusion}
\label{\detokenize{authentication:conclusion}}
\sphinxAtStartPar
Authentication ensures that only users with valid credentials can log in
to the User Access Control application.

\sphinxstepscope


\section{Authorization}
\label{\detokenize{authorization:authorization}}\label{\detokenize{authorization::doc}}

\subsection{Introduction}
\label{\detokenize{authorization:introduction}}
\sphinxAtStartPar
Authorization determines which resources an authenticated user is allowed
to access.

\sphinxAtStartPar
In the User Access Control application, each user has specific permissions
to access different pages.


\subsection{Page\sphinxhyphen{}Level Access Control}
\label{\detokenize{authorization:page-level-access-control}}
\sphinxAtStartPar
The application uses user\sphinxhyphen{}specific page permissions.

\sphinxAtStartPar
For example:

\begin{sphinxVerbatim}[commandchars=\\\{\}]
Username    Page Permissions
\PYGZhy{}\PYGZhy{}\PYGZhy{}\PYGZhy{}\PYGZhy{}\PYGZhy{}\PYGZhy{}\PYGZhy{}    \PYGZhy{}\PYGZhy{}\PYGZhy{}\PYGZhy{}\PYGZhy{}\PYGZhy{}\PYGZhy{}\PYGZhy{}\PYGZhy{}\PYGZhy{}\PYGZhy{}\PYGZhy{}\PYGZhy{}\PYGZhy{}\PYGZhy{}\PYGZhy{}
priti       page1, page2
yogesh      page3
amit        page4
admin       page1, page2, page3, page4
neha        page1, page2, page3
rahul       page2, page4
\end{sphinxVerbatim}


\subsection{Permission Checking}
\label{\detokenize{authorization:permission-checking}}
\sphinxAtStartPar
When a user requests a protected page, the application performs the following
checks:
\begin{enumerate}
\sphinxsetlistlabels{\arabic}{enumi}{enumii}{}{.}%
\item {} 
\sphinxAtStartPar
Verify whether the user is logged in.

\item {} 
\sphinxAtStartPar
Retrieve the current user’s information.

\item {} 
\sphinxAtStartPar
Retrieve the pages assigned to that user.

\item {} 
\sphinxAtStartPar
Check whether the requested page exists in the user’s permissions.

\item {} 
\sphinxAtStartPar
Allow access if permission is available.

\item {} 
\sphinxAtStartPar
Otherwise, deny access.

\end{enumerate}


\subsection{Permission Check Flow}
\label{\detokenize{authorization:permission-check-flow}}
\sphinxAtStartPar
The following diagram illustrates how the application checks user permissions
before displaying a protected page.


\subsection{Code Example}
\label{\detokenize{authorization:code-example}}
\sphinxAtStartPar
The following example demonstrates how a permission\sphinxhyphen{}checking function works.

\begin{sphinxVerbatim}[commandchars=\\\{\}]
\PYG{n}{permissions} \PYG{o}{=} \PYG{p}{\PYGZob{}}
    \PYG{l+s+s2}{\PYGZdq{}}\PYG{l+s+s2}{priti}\PYG{l+s+s2}{\PYGZdq{}}\PYG{p}{:} \PYG{p}{[}\PYG{l+s+s2}{\PYGZdq{}}\PYG{l+s+s2}{page1}\PYG{l+s+s2}{\PYGZdq{}}\PYG{p}{,} \PYG{l+s+s2}{\PYGZdq{}}\PYG{l+s+s2}{page2}\PYG{l+s+s2}{\PYGZdq{}}\PYG{p}{]}\PYG{p}{,}
    \PYG{l+s+s2}{\PYGZdq{}}\PYG{l+s+s2}{admin}\PYG{l+s+s2}{\PYGZdq{}}\PYG{p}{:} \PYG{p}{[}\PYG{l+s+s2}{\PYGZdq{}}\PYG{l+s+s2}{page1}\PYG{l+s+s2}{\PYGZdq{}}\PYG{p}{,} \PYG{l+s+s2}{\PYGZdq{}}\PYG{l+s+s2}{page2}\PYG{l+s+s2}{\PYGZdq{}}\PYG{p}{,} \PYG{l+s+s2}{\PYGZdq{}}\PYG{l+s+s2}{page3}\PYG{l+s+s2}{\PYGZdq{}}\PYG{p}{,} \PYG{l+s+s2}{\PYGZdq{}}\PYG{l+s+s2}{page4}\PYG{l+s+s2}{\PYGZdq{}}\PYG{p}{]}\PYG{p}{,}
\PYG{p}{\PYGZcb{}}

\PYG{k}{def}\PYG{+w}{ }\PYG{n+nf}{can\PYGZus{}access}\PYG{p}{(}\PYG{n}{username}\PYG{p}{,} \PYG{n}{page}\PYG{p}{)}\PYG{p}{:}
    \PYG{k}{return} \PYG{n}{page} \PYG{o+ow}{in} \PYG{n}{permissions}\PYG{o}{.}\PYG{n}{get}\PYG{p}{(}\PYG{n}{username}\PYG{p}{,} \PYG{p}{[}\PYG{p}{]}\PYG{p}{)}

\PYG{c+c1}{\PYGZsh{} Authorized access}
\PYG{n}{can\PYGZus{}access}\PYG{p}{(}\PYG{l+s+s2}{\PYGZdq{}}\PYG{l+s+s2}{priti}\PYG{l+s+s2}{\PYGZdq{}}\PYG{p}{,} \PYG{l+s+s2}{\PYGZdq{}}\PYG{l+s+s2}{page1}\PYG{l+s+s2}{\PYGZdq{}}\PYG{p}{)}  \PYG{c+c1}{\PYGZsh{} True}

\PYG{c+c1}{\PYGZsh{} Unauthorized access}
\PYG{n}{can\PYGZus{}access}\PYG{p}{(}\PYG{l+s+s2}{\PYGZdq{}}\PYG{l+s+s2}{priti}\PYG{l+s+s2}{\PYGZdq{}}\PYG{p}{,} \PYG{l+s+s2}{\PYGZdq{}}\PYG{l+s+s2}{page3}\PYG{l+s+s2}{\PYGZdq{}}\PYG{p}{)}  \PYG{c+c1}{\PYGZsh{} False}
\end{sphinxVerbatim}


\subsection{Access Control Examples}
\label{\detokenize{authorization:access-control-examples}}

\subsubsection{Example 1: Authorized Access}
\label{\detokenize{authorization:example-1-authorized-access}}
\sphinxAtStartPar
If Priti logs in and requests Page 1, access is granted because Page 1 is
included in her permissions.


\subsubsection{Example 2: Unauthorized Access}
\label{\detokenize{authorization:example-2-unauthorized-access}}
\sphinxAtStartPar
If Priti requests Page 3, access is denied because she does not have
permission to access that page.


\subsubsection{Example 3: Administrator Access}
\label{\detokenize{authorization:example-3-administrator-access}}
\sphinxAtStartPar
The administrator has permission to access all four protected pages.


\subsection{Access Denied}
\label{\detokenize{authorization:access-denied}}
\sphinxAtStartPar
When a user attempts to access a page without the required permission,
the application displays the access denied page.


\subsection{Benefits}
\label{\detokenize{authorization:benefits}}\begin{itemize}
\item {} 
\sphinxAtStartPar
Prevents unauthorized page access.

\item {} 
\sphinxAtStartPar
Provides user\sphinxhyphen{}specific permissions.

\item {} 
\sphinxAtStartPar
Separates authentication from authorization.

\item {} 
\sphinxAtStartPar
Makes access management easier.

\end{itemize}


\subsection{Conclusion}
\label{\detokenize{authorization:conclusion}}
\sphinxAtStartPar
Authorization ensures that authenticated users can access only the pages
permitted to them.

\sphinxstepscope


\section{Python Modules Documentation}
\label{\detokenize{modules:python-modules-documentation}}\label{\detokenize{modules::doc}}

\subsection{Introduction}
\label{\detokenize{modules:introduction}}
\sphinxAtStartPar
The application is divided into multiple Python modules. Each module
handles a specific responsibility to keep the code organized and maintainable.


\subsection{Module Dependency Diagram}
\label{\detokenize{modules:module-dependency-diagram}}

\subsection{main.py}
\label{\detokenize{modules:main-py}}
\sphinxAtStartPar
The main.py module is the entry point of the FastAPI application.

\sphinxAtStartPar
Responsibilities:
\begin{itemize}
\item {} 
\sphinxAtStartPar
Create the FastAPI application instance.

\item {} 
\sphinxAtStartPar
Initialize the database.

\item {} 
\sphinxAtStartPar
Configure SessionMiddleware.

\item {} 
\sphinxAtStartPar
Mount static files.

\item {} 
\sphinxAtStartPar
Register authentication and page routers.

\item {} 
\sphinxAtStartPar
Configure application exception handling.

\end{itemize}

\sphinxAtStartPar
\sphinxstylestrong{Example usage}

\begin{sphinxVerbatim}[commandchars=\\\{\}]
\PYG{c+c1}{\PYGZsh{} Run the application from the project root}
\PYG{c+c1}{\PYGZsh{}   uvicorn app.main:app \PYGZhy{}\PYGZhy{}reload}

\PYG{k+kn}{from}\PYG{+w}{ }\PYG{n+nn}{fastapi}\PYG{+w}{ }\PYG{k+kn}{import} \PYG{n}{FastAPI}
\PYG{k+kn}{from}\PYG{+w}{ }\PYG{n+nn}{starlette}\PYG{n+nn}{.}\PYG{n+nn}{middleware}\PYG{n+nn}{.}\PYG{n+nn}{sessions}\PYG{+w}{ }\PYG{k+kn}{import} \PYG{n}{SessionMiddleware}

\PYG{n}{app} \PYG{o}{=} \PYG{n}{FastAPI}\PYG{p}{(}\PYG{p}{)}
\PYG{n}{app}\PYG{o}{.}\PYG{n}{add\PYGZus{}middleware}\PYG{p}{(}\PYG{n}{SessionMiddleware}\PYG{p}{,} \PYG{n}{secret\PYGZus{}key}\PYG{o}{=}\PYG{l+s+s2}{\PYGZdq{}}\PYG{l+s+s2}{change\PYGZhy{}this\PYGZhy{}secret}\PYG{l+s+s2}{\PYGZdq{}}\PYG{p}{)}
\end{sphinxVerbatim}

\sphinxAtStartPar
\sphinxstylestrong{API reference}
\index{module@\spxentry{module}!app.main@\spxentry{app.main}}\index{app.main@\spxentry{app.main}!module@\spxentry{module}}\phantomsection\label{\detokenize{modules:module-app.main}}
\sphinxAtStartPar
Main module for the User Access Control application.

\sphinxAtStartPar
This module initializes and configures the FastAPI application,
loads environment variables, sets up session management, serves
static files, and registers application routers.
\begin{description}
\sphinxlineitem{Application Setup:}\begin{enumerate}
\sphinxsetlistlabels{\arabic}{enumi}{enumii}{}{.}%
\item {} 
\sphinxAtStartPar
Loads environment variables from the .env file using load\_dotenv().

\item {} 
\sphinxAtStartPar
Imports authentication and access\sphinxhyphen{}control routers.

\item {} 
\sphinxAtStartPar
Creates a FastAPI application instance.

\item {} 
\sphinxAtStartPar
Adds SessionMiddleware to manage user sessions securely.

\item {} 
\sphinxAtStartPar
Mounts the static directory to serve CSS, JavaScript, and other
static files.

\item {} 
\sphinxAtStartPar
Registers the authentication and access\sphinxhyphen{}control routers.

\end{enumerate}

\sphinxlineitem{Environment Variables:}\begin{description}
\sphinxlineitem{SECRET\_KEY (str): Secret key used by SessionMiddleware to sign}
\sphinxAtStartPar
session cookies. Uses a development default if not configured.

\end{description}

\sphinxlineitem{Modules:}
\sphinxAtStartPar
auth: Handles login, logout, and home page functionality.
access: Handles page1\sphinxhyphen{}page4 routes and user access permissions.
utils: Provides BASE\_DIR for locating application resources.

\sphinxlineitem{Static Files:}
\sphinxAtStartPar
URL: /static
Directory: BASE\_DIR/static

\sphinxlineitem{Application:}
\sphinxAtStartPar
app: The main FastAPI application instance.

\end{description}


\subsection{auth.py}
\label{\detokenize{modules:auth-py}}
\sphinxAtStartPar
The auth.py module handles authentication\sphinxhyphen{}related routes.

\sphinxAtStartPar
Responsibilities:
\begin{itemize}
\item {} 
\sphinxAtStartPar
Process user login requests.

\item {} 
\sphinxAtStartPar
Verify user credentials.

\item {} 
\sphinxAtStartPar
Create user sessions after successful login.

\item {} 
\sphinxAtStartPar
Handle logout requests.

\item {} 
\sphinxAtStartPar
Clear user sessions.

\end{itemize}

\sphinxAtStartPar
\sphinxstylestrong{Example usage}

\begin{sphinxVerbatim}[commandchars=\\\{\}]
\PYG{c+c1}{\PYGZsh{} POST /login  (form fields: username, password)}
\PYG{c+c1}{\PYGZsh{}}
\PYG{c+c1}{\PYGZsh{} Valid credentials   \PYGZhy{}\PYGZgt{} user ID stored in session, redirect to home page}
\PYG{c+c1}{\PYGZsh{} Invalid credentials \PYGZhy{}\PYGZgt{} login.html shown again with an error message}
\PYG{c+c1}{\PYGZsh{}}
\PYG{c+c1}{\PYGZsh{} GET /logout}
\PYG{c+c1}{\PYGZsh{}}
\PYG{c+c1}{\PYGZsh{} Session cleared \PYGZhy{}\PYGZgt{} redirect to login page}
\end{sphinxVerbatim}

\sphinxAtStartPar
\sphinxstylestrong{API reference}
\index{module@\spxentry{module}!app.auth@\spxentry{app.auth}}\index{app.auth@\spxentry{app.auth}!module@\spxentry{module}}\phantomsection\label{\detokenize{modules:module-app.auth}}
\sphinxAtStartPar
Authentication module for the User Access Control application.

\sphinxAtStartPar
This module defines the authentication API router and imports the
functions required to handle user login, logout, session management,
and home page functionality.

\sphinxAtStartPar
It uses FastAPI form handling, database queries, password verification,
and session\sphinxhyphen{}based authentication to manage user access.
\begin{description}
\sphinxlineitem{Imports:}
\sphinxAtStartPar
escape: Escapes special HTML characters to help prevent HTML injection.
APIRouter: Creates a router for registering authentication endpoints.
Request: Provides access to request data and session information.
Form: Extracts form data submitted by users.
HTMLResponse: Returns HTML content as an HTTP response.
RedirectResponse: Redirects users to another URL.
Response: Base response class for HTML and redirect responses.
get\_user\_by\_username: Retrieves a user record using their username.
verify\_password: Checks the entered password against the stored hash.
load\_html: Loads HTML templates from the application’s HTML directory.
current\_user: Retrieves the currently logged\sphinxhyphen{}in user’s database record.
get\_user\_pages: Retrieves the pages assigned to a user.

\sphinxlineitem{Global Variables:}
\sphinxAtStartPar
router (APIRouter): Router used to register authentication endpoints.

\sphinxlineitem{Authentication Features:}\begin{itemize}
\item {} 
\sphinxAtStartPar
User login and password verification.

\item {} 
\sphinxAtStartPar
Session management after successful login.

\item {} 
\sphinxAtStartPar
Logout and session clearing.

\item {} 
\sphinxAtStartPar
Home page access for authenticated users.

\item {} 
\sphinxAtStartPar
Retrieving user\sphinxhyphen{}specific page permissions.

\end{itemize}

\end{description}
\index{home() (in module app.auth)@\spxentry{home()}\spxextra{in module app.auth}}

\begin{fulllineitems}
\phantomsection\label{\detokenize{modules:app.auth.home}}
\pysigstartsignatures
\pysiglinewithargsret
{\sphinxcode{\sphinxupquote{app.auth.}}\sphinxbfcode{\sphinxupquote{home}}}
{\sphinxparam{\DUrole{n}{request}\DUrole{p}{:}\DUrole{w}{ }\DUrole{n}{Request}}}
{{ $\rightarrow$ Response}}
\pysigstopsignatures
\sphinxAtStartPar
Display the home page with the logged\sphinxhyphen{}in user’s details
and assigned page links.
\begin{quote}\begin{description}
\sphinxlineitem{Parameters}
\sphinxAtStartPar
\sphinxstyleliteralstrong{\sphinxupquote{request}} (\sphinxstyleliteralemphasis{\sphinxupquote{Request}}) \textendash{} Incoming HTTP request object.
Used to identify the currently logged\sphinxhyphen{}in user
and access session information.

\end{description}\end{quote}
\begin{description}
\sphinxlineitem{What it does:}\begin{enumerate}
\sphinxsetlistlabels{\arabic}{enumi}{enumii}{}{.}%
\item {} 
\sphinxAtStartPar
Handles GET requests to the “/home” URL.

\item {} 
\sphinxAtStartPar
Checks whether the user is logged in using current\_user().

\item {} 
\sphinxAtStartPar
If the user is not logged in, redirects them to “/login”.

\item {} 
\sphinxAtStartPar
Retrieves the list of pages assigned to the user
using get\_user\_pages().

\item {} 
\sphinxAtStartPar
If the user has assigned pages:
\sphinxhyphen{} Generates HTML list items containing links to each page.
\sphinxhyphen{} Replaces “page” in page names with “Page “ for display.

\item {} 
\sphinxAtStartPar
If no pages are assigned, displays a message:
“No pages assigned to you yet.”

\item {} 
\sphinxAtStartPar
Loads the home.html template.

\item {} 
\sphinxAtStartPar
Replaces the “\_\_NAME\_\_” placeholder with the user’s
escaped full name.

\item {} 
\sphinxAtStartPar
Replaces the “\_\_LINKS\_\_” placeholder with the generated
page links or the no\sphinxhyphen{}pages message.

\item {} 
\sphinxAtStartPar
Returns the updated home page as an HTML response.

\end{enumerate}

\end{description}
\begin{quote}\begin{description}
\sphinxlineitem{Returns}
\sphinxAtStartPar
\begin{description}
\sphinxlineitem{RedirectResponse: Redirects unauthenticated users}
\sphinxAtStartPar
to “/login” with status code 303.

\sphinxlineitem{HTMLResponse: Displays the home page with the user’s}
\sphinxAtStartPar
name and assigned page links.

\end{description}


\sphinxlineitem{Return type}
\sphinxAtStartPar
Response

\end{description}\end{quote}
\subsubsection*{Example}
\begin{description}
\sphinxlineitem{Logged\sphinxhyphen{}in user with assigned pages:}
\sphinxAtStartPar
Displays the home page with the user’s name
and available page links.

\sphinxlineitem{Logged\sphinxhyphen{}in user without assigned pages:}
\sphinxAtStartPar
Displays “No pages assigned to you yet.”

\sphinxlineitem{Unauthenticated user:}
\sphinxAtStartPar
Redirects to “/login”.

\end{description}

\end{fulllineitems}

\index{login() (in module app.auth)@\spxentry{login()}\spxextra{in module app.auth}}

\begin{fulllineitems}
\phantomsection\label{\detokenize{modules:app.auth.login}}
\pysigstartsignatures
\pysiglinewithargsret
{\sphinxcode{\sphinxupquote{app.auth.}}\sphinxbfcode{\sphinxupquote{login}}}
{\sphinxparam{\DUrole{n}{request}\DUrole{p}{:}\DUrole{w}{ }\DUrole{n}{Request}}\sphinxparamcomma \sphinxparam{\DUrole{n}{username}\DUrole{p}{:}\DUrole{w}{ }\DUrole{n}{str}\DUrole{w}{ }\DUrole{o}{=}\DUrole{w}{ }\DUrole{default_value}{Form(PydanticUndefined)}}\sphinxparamcomma \sphinxparam{\DUrole{n}{password}\DUrole{p}{:}\DUrole{w}{ }\DUrole{n}{str}\DUrole{w}{ }\DUrole{o}{=}\DUrole{w}{ }\DUrole{default_value}{Form(PydanticUndefined)}}}
{{ $\rightarrow$ Response}}
\pysigstopsignatures
\sphinxAtStartPar
Authenticate a user using their username and password.
\begin{quote}\begin{description}
\sphinxlineitem{Parameters}\begin{itemize}
\item {} 
\sphinxAtStartPar
\sphinxstyleliteralstrong{\sphinxupquote{request}} (\sphinxstyleliteralemphasis{\sphinxupquote{Request}}) \textendash{} Incoming HTTP request object.
Used to access and manage the user session.

\item {} 
\sphinxAtStartPar
\sphinxstyleliteralstrong{\sphinxupquote{username}} (\sphinxstyleliteralemphasis{\sphinxupquote{str}}) \textendash{} Username entered in the login form.
Required form field.

\item {} 
\sphinxAtStartPar
\sphinxstyleliteralstrong{\sphinxupquote{password}} (\sphinxstyleliteralemphasis{\sphinxupquote{str}}) \textendash{} Password entered in the login form.
Required form field.

\end{itemize}

\end{description}\end{quote}
\begin{description}
\sphinxlineitem{What it does:}\begin{enumerate}
\sphinxsetlistlabels{\arabic}{enumi}{enumii}{}{.}%
\item {} 
\sphinxAtStartPar
Receives the username and password submitted through
the login form.

\item {} 
\sphinxAtStartPar
Removes leading and trailing whitespace from the username.

\item {} 
\sphinxAtStartPar
Fetches the user details from the database using
get\_user\_by\_username().

\item {} 
\sphinxAtStartPar
Checks whether the user exists.

\item {} 
\sphinxAtStartPar
Verifies the entered password against the stored
password hash using verify\_password().

\item {} 
\sphinxAtStartPar
If the username or password is incorrect:
\sphinxhyphen{} Displays an error message.
\sphinxhyphen{} Returns the login page with HTTP status code 401.

\item {} 
\sphinxAtStartPar
If authentication is successful:
\sphinxhyphen{} Clears the existing session.
\sphinxhyphen{} Stores the user’s ID in the session.
\sphinxhyphen{} Redirects the user to the home page.

\end{enumerate}

\end{description}
\begin{quote}\begin{description}
\sphinxlineitem{Returns}
\sphinxAtStartPar
\begin{description}
\sphinxlineitem{HTMLResponse: Returns the login page with an error}
\sphinxAtStartPar
message if authentication fails.

\sphinxlineitem{RedirectResponse: Redirects the authenticated user}
\sphinxAtStartPar
to “/home” with HTTP status code 303.

\end{description}


\sphinxlineitem{Return type}
\sphinxAtStartPar
HTMLResponse | RedirectResponse

\end{description}\end{quote}
\subsubsection*{Example}
\begin{description}
\sphinxlineitem{Input:}
\sphinxAtStartPar
username = “admin”
password = “admin123”

\sphinxlineitem{Successful login:}
\sphinxAtStartPar
Redirects to “/home”.

\sphinxlineitem{Failed login:}
\sphinxAtStartPar
Displays “Wrong username or password.”
and returns HTTP 401.

\end{description}

\end{fulllineitems}

\index{login\_form() (in module app.auth)@\spxentry{login\_form()}\spxextra{in module app.auth}}

\begin{fulllineitems}
\phantomsection\label{\detokenize{modules:app.auth.login_form}}
\pysigstartsignatures
\pysiglinewithargsret
{\sphinxcode{\sphinxupquote{app.auth.}}\sphinxbfcode{\sphinxupquote{login\_form}}}
{}
{{ $\rightarrow$ HTMLResponse}}
\pysigstopsignatures
\sphinxAtStartPar
Display the login page to the user.
\begin{quote}\begin{description}
\sphinxlineitem{Parameters}
\sphinxAtStartPar
\sphinxstyleliteralstrong{\sphinxupquote{None}} \textendash{} This function does not require any input parameters.

\end{description}\end{quote}
\begin{description}
\sphinxlineitem{What it does:}\begin{enumerate}
\sphinxsetlistlabels{\arabic}{enumi}{enumii}{}{.}%
\item {} 
\sphinxAtStartPar
Handles GET requests to the “/login” URL.

\item {} 
\sphinxAtStartPar
Loads the login.html file.

\item {} 
\sphinxAtStartPar
Replaces the “\_\_ERROR\_\_” placeholder with an empty string.
This ensures that no error message is displayed initially.

\item {} 
\sphinxAtStartPar
Returns the updated HTML content as a response.

\end{enumerate}

\end{description}
\begin{quote}\begin{description}
\sphinxlineitem{Returns}
\sphinxAtStartPar
\begin{description}
\sphinxlineitem{Returns the login page HTML content}
\sphinxAtStartPar
with the error message placeholder cleared.

\end{description}


\sphinxlineitem{Return type}
\sphinxAtStartPar
HTMLResponse

\end{description}\end{quote}
\subsubsection*{Example}
\begin{description}
\sphinxlineitem{Request:}
\sphinxAtStartPar
GET /login

\sphinxlineitem{Output:}
\sphinxAtStartPar
Displays the login page without an error message.

\end{description}

\end{fulllineitems}

\index{logout() (in module app.auth)@\spxentry{logout()}\spxextra{in module app.auth}}

\begin{fulllineitems}
\phantomsection\label{\detokenize{modules:app.auth.logout}}
\pysigstartsignatures
\pysiglinewithargsret
{\sphinxcode{\sphinxupquote{app.auth.}}\sphinxbfcode{\sphinxupquote{logout}}}
{\sphinxparam{\DUrole{n}{request}\DUrole{p}{:}\DUrole{w}{ }\DUrole{n}{Request}}}
{{ $\rightarrow$ RedirectResponse}}
\pysigstopsignatures
\sphinxAtStartPar
Log out the currently authenticated user.
\begin{quote}\begin{description}
\sphinxlineitem{Parameters}
\sphinxAtStartPar
\sphinxstyleliteralstrong{\sphinxupquote{request}} (\sphinxstyleliteralemphasis{\sphinxupquote{Request}}) \textendash{} Incoming HTTP request object.
Used to access and manage the user’s session.

\end{description}\end{quote}
\begin{description}
\sphinxlineitem{What it does:}\begin{enumerate}
\sphinxsetlistlabels{\arabic}{enumi}{enumii}{}{.}%
\item {} 
\sphinxAtStartPar
Handles GET requests to the “/logout” URL.

\item {} 
\sphinxAtStartPar
Clears all data stored in the user’s session.

\item {} 
\sphinxAtStartPar
Removes the user’s login information from the session.

\item {} 
\sphinxAtStartPar
Redirects the user to the login page.

\item {} 
\sphinxAtStartPar
Uses HTTP status code 303 for redirection.

\end{enumerate}

\end{description}
\begin{quote}\begin{description}
\sphinxlineitem{Returns}
\sphinxAtStartPar
\begin{description}
\sphinxlineitem{Redirects the user to “/login”}
\sphinxAtStartPar
after clearing the session.

\end{description}


\sphinxlineitem{Return type}
\sphinxAtStartPar
RedirectResponse

\end{description}\end{quote}
\subsubsection*{Example}
\begin{description}
\sphinxlineitem{Request:}
\sphinxAtStartPar
GET /logout

\sphinxlineitem{Output:}
\sphinxAtStartPar
Clears the user session and redirects to “/login”.

\end{description}

\end{fulllineitems}

\index{root() (in module app.auth)@\spxentry{root()}\spxextra{in module app.auth}}

\begin{fulllineitems}
\phantomsection\label{\detokenize{modules:app.auth.root}}
\pysigstartsignatures
\pysiglinewithargsret
{\sphinxcode{\sphinxupquote{app.auth.}}\sphinxbfcode{\sphinxupquote{root}}}
{\sphinxparam{\DUrole{n}{request}\DUrole{p}{:}\DUrole{w}{ }\DUrole{n}{Request}}}
{{ $\rightarrow$ RedirectResponse}}
\pysigstopsignatures
\sphinxAtStartPar
Handle requests to the root URL (“/”) and redirect users
based on their login status.
\begin{quote}\begin{description}
\sphinxlineitem{Parameters}
\sphinxAtStartPar
\sphinxstyleliteralstrong{\sphinxupquote{request}} (\sphinxstyleliteralemphasis{\sphinxupquote{Request}}) \textendash{} Incoming HTTP request object.
Contains the user’s session information.

\end{description}\end{quote}
\begin{description}
\sphinxlineitem{What it does:}\begin{enumerate}
\sphinxsetlistlabels{\arabic}{enumi}{enumii}{}{.}%
\item {} 
\sphinxAtStartPar
Receives a request when the user visits the root URL.

\item {} 
\sphinxAtStartPar
Checks whether the user’s session contains a “user\_id”.

\item {} 
\sphinxAtStartPar
If user\_id exists, redirects the user to “/home”.

\item {} 
\sphinxAtStartPar
If user\_id does not exist, redirects the user to “/login”.

\item {} 
\sphinxAtStartPar
Uses HTTP status code 303 to perform the redirect.

\end{enumerate}

\end{description}
\begin{quote}\begin{description}
\sphinxlineitem{Returns}
\sphinxAtStartPar
\begin{description}
\sphinxlineitem{Redirects the user to either:}\begin{itemize}
\item {} 
\sphinxAtStartPar
”/home” if the user is logged in.

\item {} 
\sphinxAtStartPar
”/login” if the user is not logged in.

\end{itemize}

\end{description}


\sphinxlineitem{Return type}
\sphinxAtStartPar
RedirectResponse

\end{description}\end{quote}
\subsubsection*{Example}
\begin{description}
\sphinxlineitem{Logged\sphinxhyphen{}in user:}
\sphinxAtStartPar
GET /
Redirects to /home

\sphinxlineitem{Not logged\sphinxhyphen{}in user:}
\sphinxAtStartPar
GET /
Redirects to /login

\end{description}

\end{fulllineitems}



\subsection{access.py}
\label{\detokenize{modules:access-py}}
\sphinxAtStartPar
The access.py module handles page\sphinxhyphen{}level authorization.

\sphinxAtStartPar
Responsibilities:
\begin{itemize}
\item {} 
\sphinxAtStartPar
Define or retrieve user page permissions.

\item {} 
\sphinxAtStartPar
Check whether a user can access a requested page.

\item {} 
\sphinxAtStartPar
Protect restricted pages.

\item {} 
\sphinxAtStartPar
Handle unauthorized access.

\end{itemize}

\sphinxAtStartPar
\sphinxstylestrong{Example usage}

\begin{sphinxVerbatim}[commandchars=\\\{\}]
\PYG{k+kn}{from}\PYG{+w}{ }\PYG{n+nn}{app}\PYG{n+nn}{.}\PYG{n+nn}{access}\PYG{+w}{ }\PYG{k+kn}{import} \PYG{n}{check\PYGZus{}page\PYGZus{}access}

\PYG{k}{if} \PYG{n}{check\PYGZus{}page\PYGZus{}access}\PYG{p}{(}\PYG{n}{username}\PYG{o}{=}\PYG{l+s+s2}{\PYGZdq{}}\PYG{l+s+s2}{priti}\PYG{l+s+s2}{\PYGZdq{}}\PYG{p}{,} \PYG{n}{page}\PYG{o}{=}\PYG{l+s+s2}{\PYGZdq{}}\PYG{l+s+s2}{page1}\PYG{l+s+s2}{\PYGZdq{}}\PYG{p}{)}\PYG{p}{:}
    \PYG{n+nb}{print}\PYG{p}{(}\PYG{l+s+s2}{\PYGZdq{}}\PYG{l+s+s2}{Show page1.html}\PYG{l+s+s2}{\PYGZdq{}}\PYG{p}{)}      \PYG{c+c1}{\PYGZsh{} Priti has page1 permission}
\PYG{k}{else}\PYG{p}{:}
    \PYG{n+nb}{print}\PYG{p}{(}\PYG{l+s+s2}{\PYGZdq{}}\PYG{l+s+s2}{Show denied.html}\PYG{l+s+s2}{\PYGZdq{}}\PYG{p}{)}

\PYG{n}{check\PYGZus{}page\PYGZus{}access}\PYG{p}{(}\PYG{l+s+s2}{\PYGZdq{}}\PYG{l+s+s2}{priti}\PYG{l+s+s2}{\PYGZdq{}}\PYG{p}{,} \PYG{l+s+s2}{\PYGZdq{}}\PYG{l+s+s2}{page3}\PYG{l+s+s2}{\PYGZdq{}}\PYG{p}{)}   \PYG{c+c1}{\PYGZsh{} False \PYGZhy{}\PYGZgt{} denied.html}
\PYG{n}{check\PYGZus{}page\PYGZus{}access}\PYG{p}{(}\PYG{l+s+s2}{\PYGZdq{}}\PYG{l+s+s2}{admin}\PYG{l+s+s2}{\PYGZdq{}}\PYG{p}{,} \PYG{l+s+s2}{\PYGZdq{}}\PYG{l+s+s2}{page3}\PYG{l+s+s2}{\PYGZdq{}}\PYG{p}{)}   \PYG{c+c1}{\PYGZsh{} True  \PYGZhy{}\PYGZgt{} page3.html}
\end{sphinxVerbatim}

\sphinxAtStartPar
\sphinxstylestrong{API reference}
\index{module@\spxentry{module}!app.access@\spxentry{app.access}}\index{app.access@\spxentry{app.access}!module@\spxentry{module}}\phantomsection\label{\detokenize{modules:module-app.access}}
\sphinxAtStartPar
Access Control module for the User Access Control application.

\sphinxAtStartPar
This module defines the API router and configuration required to manage
protected pages and control user access based on assigned permissions.

\sphinxAtStartPar
It uses FastAPI routing, database connections, session\sphinxhyphen{}based authentication,
and HTML response handling to provide authorized access to application pages.
\begin{description}
\sphinxlineitem{Imports:}
\sphinxAtStartPar
escape: Escapes special HTML characters to help prevent HTML injection.
APIRouter: Creates a router for registering application endpoints.
Request: Provides access to request and session information.
HTMLResponse: Returns HTML content as an HTTP response.
RedirectResponse: Redirects users to another URL.
get\_connection: Establishes a connection to the MySQL database.
load\_html: Loads HTML files from the application’s HTML directory.
current\_user: Retrieves the currently authenticated user.
Callable: Represents a function that can be called.
Response: Base response class for HTML and redirect responses.

\sphinxlineitem{Global Variables:}
\sphinxAtStartPar
router (APIRouter): Router used to register access\sphinxhyphen{}control endpoints.
PAGES (list{[}str{]}): List of protected page names managed by this module.

\sphinxlineitem{Protected Pages:}
\sphinxAtStartPar
page1, page2, page3, page4

\end{description}
\index{get\_user\_pages() (in module app.access)@\spxentry{get\_user\_pages()}\spxextra{in module app.access}}

\begin{fulllineitems}
\phantomsection\label{\detokenize{modules:app.access.get_user_pages}}
\pysigstartsignatures
\pysiglinewithargsret
{\sphinxcode{\sphinxupquote{app.access.}}\sphinxbfcode{\sphinxupquote{get\_user\_pages}}}
{\sphinxparam{\DUrole{n}{user\_id}}}
{}
\pysigstopsignatures
\sphinxAtStartPar
Get all pages that a specific user is allowed to access.
\begin{quote}\begin{description}
\sphinxlineitem{Parameters}
\sphinxAtStartPar
\sphinxstyleliteralstrong{\sphinxupquote{user\_id}} (\sphinxstyleliteralemphasis{\sphinxupquote{int}}\sphinxstyleliteralemphasis{\sphinxupquote{ or }}\sphinxstyleliteralemphasis{\sphinxupquote{str}}) \textendash{} Unique ID of the user.

\end{description}\end{quote}
\begin{description}
\sphinxlineitem{What it does:}\begin{enumerate}
\sphinxsetlistlabels{\arabic}{enumi}{enumii}{}{.}%
\item {} 
\sphinxAtStartPar
Connects to the database.

\item {} 
\sphinxAtStartPar
Fetches all page names assigned to the given user\_id
from the user\_pages table.

\item {} 
\sphinxAtStartPar
Sorts the page names alphabetically.

\item {} 
\sphinxAtStartPar
Closes the database connection.

\end{enumerate}

\end{description}
\begin{quote}\begin{description}
\sphinxlineitem{Returns}
\sphinxAtStartPar
\begin{description}
\sphinxlineitem{A list of page names assigned to the user.}
\sphinxAtStartPar
Returns an empty list if no pages are assigned.

\end{description}


\sphinxlineitem{Return type}
\sphinxAtStartPar
list{[}str{]}

\end{description}\end{quote}
\subsubsection*{Example}
\begin{description}
\sphinxlineitem{Input:}
\sphinxAtStartPar
get\_user\_pages(101)

\sphinxlineitem{Output:}
\sphinxAtStartPar
{[}‘Attendance’, ‘Dashboard’, ‘Employee’{]}

\end{description}

\end{fulllineitems}

\index{make\_page\_route() (in module app.access)@\spxentry{make\_page\_route()}\spxextra{in module app.access}}

\begin{fulllineitems}
\phantomsection\label{\detokenize{modules:app.access.make_page_route}}
\pysigstartsignatures
\pysiglinewithargsret
{\sphinxcode{\sphinxupquote{app.access.}}\sphinxbfcode{\sphinxupquote{make\_page\_route}}}
{\sphinxparam{\DUrole{n}{page\_name}\DUrole{p}{:}\DUrole{w}{ }\DUrole{n}{str}}}
{{ $\rightarrow$ Callable\DUrole{p}{{[}}\DUrole{p}{{[}}Request\DUrole{p}{{]}}\DUrole{p}{,}\DUrole{w}{ }Response\DUrole{p}{{]}}}}
\pysigstopsignatures
\sphinxAtStartPar
Create a page route handler with authentication
and authorization checks.
\begin{quote}\begin{description}
\sphinxlineitem{Parameters}
\sphinxAtStartPar
\sphinxstyleliteralstrong{\sphinxupquote{page\_name}} (\sphinxstyleliteralemphasis{\sphinxupquote{str}}) \textendash{} Name of the HTML page to access.
Example: “page1”, “page2”, “dashboard”.

\end{description}\end{quote}
\begin{description}
\sphinxlineitem{What it does:}\begin{enumerate}
\sphinxsetlistlabels{\arabic}{enumi}{enumii}{}{.}%
\item {} 
\sphinxAtStartPar
Creates an inner function named page\_view().

\item {} 
\sphinxAtStartPar
Checks whether the user is logged in.

\item {} 
\sphinxAtStartPar
Redirects unauthenticated users to the login page.

\item {} 
\sphinxAtStartPar
Checks whether the requested page is assigned
to the logged\sphinxhyphen{}in user.

\item {} 
\sphinxAtStartPar
If access is denied, displays the denied.html page.

\item {} 
\sphinxAtStartPar
If access is allowed, displays the requested HTML page.

\end{enumerate}

\end{description}
\begin{quote}\begin{description}
\sphinxlineitem{Returns}
\sphinxAtStartPar
Returns a page\_view function that accepts an HTTP
Request and returns an HTTP Response.

\sphinxlineitem{Return type}
\sphinxAtStartPar
Callable{[}{[}Request{]}, Response{]}

\end{description}\end{quote}
\subsubsection*{Example}

\sphinxAtStartPar
make\_page\_route(“page1”)

\sphinxAtStartPar
Returns a page handler for page1.

\end{fulllineitems}



\subsection{db.py}
\label{\detokenize{modules:db-py}}
\sphinxAtStartPar
The db.py module handles database operations.

\sphinxAtStartPar
Responsibilities:
\begin{itemize}
\item {} 
\sphinxAtStartPar
Connect to the MySQL database.

\item {} 
\sphinxAtStartPar
Retrieve user records.

\item {} 
\sphinxAtStartPar
Retrieve user permissions.

\item {} 
\sphinxAtStartPar
Perform database\sphinxhyphen{}related operations.

\end{itemize}

\sphinxAtStartPar
\sphinxstylestrong{Example usage}

\begin{sphinxVerbatim}[commandchars=\\\{\}]
\PYG{k+kn}{from}\PYG{+w}{ }\PYG{n+nn}{app}\PYG{n+nn}{.}\PYG{n+nn}{db}\PYG{+w}{ }\PYG{k+kn}{import} \PYG{n}{get\PYGZus{}connection}

\PYG{n}{conn} \PYG{o}{=} \PYG{n}{get\PYGZus{}connection}\PYG{p}{(}\PYG{p}{)}
\PYG{n}{cursor} \PYG{o}{=} \PYG{n}{conn}\PYG{o}{.}\PYG{n}{cursor}\PYG{p}{(}\PYG{p}{)}
\PYG{n}{cursor}\PYG{o}{.}\PYG{n}{execute}\PYG{p}{(}\PYG{l+s+s2}{\PYGZdq{}}\PYG{l+s+s2}{SELECT username FROM users WHERE username = }\PYG{l+s+si}{\PYGZpc{}s}\PYG{l+s+s2}{\PYGZdq{}}\PYG{p}{,} \PYG{p}{(}\PYG{l+s+s2}{\PYGZdq{}}\PYG{l+s+s2}{priti}\PYG{l+s+s2}{\PYGZdq{}}\PYG{p}{,}\PYG{p}{)}\PYG{p}{)}
\PYG{n+nb}{print}\PYG{p}{(}\PYG{n}{cursor}\PYG{o}{.}\PYG{n}{fetchone}\PYG{p}{(}\PYG{p}{)}\PYG{p}{)}
\PYG{n}{conn}\PYG{o}{.}\PYG{n}{close}\PYG{p}{(}\PYG{p}{)}
\end{sphinxVerbatim}

\sphinxAtStartPar
\sphinxstylestrong{API reference}
\index{module@\spxentry{module}!app.db@\spxentry{app.db}}\index{app.db@\spxentry{app.db}!module@\spxentry{module}}\phantomsection\label{\detokenize{modules:module-app.db}}
\sphinxAtStartPar
Database configuration module for the User Access Control application.

\sphinxAtStartPar
This module imports the required libraries and loads environment
variables used for connecting to the MySQL database.
\begin{description}
\sphinxlineitem{Imports:}\begin{description}
\sphinxlineitem{os: Provides functions to access environment variables and}
\sphinxAtStartPar
handle operating system paths.

\end{description}

\sphinxAtStartPar
pymysql: Provides MySQL database connectivity and query execution.
load\_dotenv: Loads database configuration variables from a .env file.

\sphinxlineitem{Configuration:}
\sphinxAtStartPar
load\_dotenv(): Loads environment variables from the .env file
into the application’s environment.

\sphinxlineitem{Purpose:}
\sphinxAtStartPar
Provides the database configuration required by other modules
to establish MySQL database connections.

\end{description}
\index{get\_connection() (in module app.db)@\spxentry{get\_connection()}\spxextra{in module app.db}}

\begin{fulllineitems}
\phantomsection\label{\detokenize{modules:app.db.get_connection}}
\pysigstartsignatures
\pysiglinewithargsret
{\sphinxcode{\sphinxupquote{app.db.}}\sphinxbfcode{\sphinxupquote{get\_connection}}}
{}
{}
\pysigstopsignatures
\sphinxAtStartPar
Create and return a connection to the MySQL database.
\begin{description}
\sphinxlineitem{Reads the database configuration from environment variables:}
\sphinxAtStartPar
DB\_HOST: Database server hostname (default: localhost).
DB\_PORT: Database server port (default: 3306).
DB\_USER: Database username (default: root).
DB\_PASSWORD: Database password.
DB\_NAME: Database name (default: access\_demo).

\end{description}

\sphinxAtStartPar
Uses DictCursor so query results are returned as dictionaries,
allowing columns to be accessed by their names.
\begin{quote}\begin{description}
\sphinxlineitem{Returns}
\sphinxAtStartPar
An active MySQL database connection.

\sphinxlineitem{Return type}
\sphinxAtStartPar
pymysql.connections.Connection

\sphinxlineitem{Raises}\begin{itemize}
\item {} 
\sphinxAtStartPar
\sphinxstyleliteralstrong{\sphinxupquote{pymysql.MySQLError}} \textendash{} If the database connection cannot be established.

\item {} 
\sphinxAtStartPar
\sphinxstyleliteralstrong{\sphinxupquote{ValueError}} \textendash{} If DB\_PORT contains an invalid integer value.

\end{itemize}

\end{description}\end{quote}
\subsubsection*{Example}

\sphinxAtStartPar
conn = get\_connection()
cursor = conn.cursor()
cursor.execute(“SELECT * FROM users”)
users = cursor.fetchall()
conn.close()

\end{fulllineitems}

\index{get\_user\_by\_id() (in module app.db)@\spxentry{get\_user\_by\_id()}\spxextra{in module app.db}}

\begin{fulllineitems}
\phantomsection\label{\detokenize{modules:app.db.get_user_by_id}}
\pysigstartsignatures
\pysiglinewithargsret
{\sphinxcode{\sphinxupquote{app.db.}}\sphinxbfcode{\sphinxupquote{get\_user\_by\_id}}}
{\sphinxparam{\DUrole{n}{user\_id}}}
{}
\pysigstopsignatures
\sphinxAtStartPar
Retrieve a user record from the database using their user ID.

\sphinxAtStartPar
Establishes a database connection, executes a parameterized SQL query
to find the user with the specified ID, and fetches the matching record.
\begin{quote}\begin{description}
\sphinxlineitem{Parameters}
\sphinxAtStartPar
\sphinxstyleliteralstrong{\sphinxupquote{user\_id}} (\sphinxstyleliteralemphasis{\sphinxupquote{int}}) \textendash{} The unique ID of the user to retrieve.

\sphinxlineitem{Returns}
\sphinxAtStartPar
A dictionary containing the user’s database record
if a matching user is found; otherwise, None.

\sphinxlineitem{Return type}
\sphinxAtStartPar
dict | None

\sphinxlineitem{Raises}
\sphinxAtStartPar
\sphinxstyleliteralstrong{\sphinxupquote{pymysql.MySQLError}} \textendash{} If a database operation fails.

\end{description}\end{quote}
\subsubsection*{Example}

\sphinxAtStartPar
user = get\_user\_by\_id(1)
\begin{description}
\sphinxlineitem{if user:}
\sphinxAtStartPar
print(user{[}“username”{]})

\sphinxlineitem{else:}
\sphinxAtStartPar
print(“User not found.”)

\end{description}

\end{fulllineitems}

\index{get\_user\_by\_username() (in module app.db)@\spxentry{get\_user\_by\_username()}\spxextra{in module app.db}}

\begin{fulllineitems}
\phantomsection\label{\detokenize{modules:app.db.get_user_by_username}}
\pysigstartsignatures
\pysiglinewithargsret
{\sphinxcode{\sphinxupquote{app.db.}}\sphinxbfcode{\sphinxupquote{get\_user\_by\_username}}}
{\sphinxparam{\DUrole{n}{username}}}
{}
\pysigstopsignatures
\sphinxAtStartPar
Retrieve a user record from the database using their username.

\sphinxAtStartPar
Establishes a database connection, executes a parameterized SQL query
to find the matching user, and fetches the first matching record.
\begin{quote}\begin{description}
\sphinxlineitem{Parameters}
\sphinxAtStartPar
\sphinxstyleliteralstrong{\sphinxupquote{username}} (\sphinxstyleliteralemphasis{\sphinxupquote{str}}) \textendash{} The username of the user to search for.

\sphinxlineitem{Returns}
\sphinxAtStartPar
A dictionary containing the user’s database record
if a matching user is found; otherwise, None.

\sphinxlineitem{Return type}
\sphinxAtStartPar
dict | None

\sphinxlineitem{Raises}
\sphinxAtStartPar
\sphinxstyleliteralstrong{\sphinxupquote{pymysql.MySQLError}} \textendash{} If a database operation fails.

\end{description}\end{quote}
\subsubsection*{Example}

\sphinxAtStartPar
user = get\_user\_by\_username(“admin”)
\begin{description}
\sphinxlineitem{if user:}
\sphinxAtStartPar
print(user{[}“username”{]})

\sphinxlineitem{else:}
\sphinxAtStartPar
print(“User not found.”)

\end{description}

\end{fulllineitems}



\subsection{security.py}
\label{\detokenize{modules:security-py}}
\sphinxAtStartPar
The security.py module handles password security.

\sphinxAtStartPar
Responsibilities:
\begin{itemize}
\item {} 
\sphinxAtStartPar
Generate password hashes using bcrypt.

\item {} 
\sphinxAtStartPar
Verify passwords against stored hashes.

\end{itemize}

\sphinxAtStartPar
\sphinxstylestrong{Example usage}

\begin{sphinxVerbatim}[commandchars=\\\{\}]
\PYG{k+kn}{from}\PYG{+w}{ }\PYG{n+nn}{app}\PYG{n+nn}{.}\PYG{n+nn}{security}\PYG{+w}{ }\PYG{k+kn}{import} \PYG{n}{hash\PYGZus{}password}\PYG{p}{,} \PYG{n}{verify\PYGZus{}password}

\PYG{n}{hashed} \PYG{o}{=} \PYG{n}{hash\PYGZus{}password}\PYG{p}{(}\PYG{l+s+s2}{\PYGZdq{}}\PYG{l+s+s2}{MySecret@123}\PYG{l+s+s2}{\PYGZdq{}}\PYG{p}{)}
\PYG{n+nb}{print}\PYG{p}{(}\PYG{n}{verify\PYGZus{}password}\PYG{p}{(}\PYG{l+s+s2}{\PYGZdq{}}\PYG{l+s+s2}{MySecret@123}\PYG{l+s+s2}{\PYGZdq{}}\PYG{p}{,} \PYG{n}{hashed}\PYG{p}{)}\PYG{p}{)}   \PYG{c+c1}{\PYGZsh{} True}
\PYG{n+nb}{print}\PYG{p}{(}\PYG{n}{verify\PYGZus{}password}\PYG{p}{(}\PYG{l+s+s2}{\PYGZdq{}}\PYG{l+s+s2}{wrong}\PYG{l+s+s2}{\PYGZdq{}}\PYG{p}{,} \PYG{n}{hashed}\PYG{p}{)}\PYG{p}{)}          \PYG{c+c1}{\PYGZsh{} False}
\end{sphinxVerbatim}

\sphinxAtStartPar
\sphinxstylestrong{API reference}
\index{module@\spxentry{module}!app.security@\spxentry{app.security}}\index{app.security@\spxentry{app.security}!module@\spxentry{module}}\index{hash\_password() (in module app.security)@\spxentry{hash\_password()}\spxextra{in module app.security}}\phantomsection\label{\detokenize{modules:module-app.security}}

\begin{fulllineitems}
\phantomsection\label{\detokenize{modules:app.security.hash_password}}
\pysigstartsignatures
\pysiglinewithargsret
{\sphinxcode{\sphinxupquote{app.security.}}\sphinxbfcode{\sphinxupquote{hash\_password}}}
{\sphinxparam{\DUrole{n}{plain\_password}\DUrole{p}{:}\DUrole{w}{ }\DUrole{n}{str}}}
{{ $\rightarrow$ str}}
\pysigstopsignatures
\sphinxAtStartPar
Hash a plain\sphinxhyphen{}text password using the bcrypt hashing algorithm.

\sphinxAtStartPar
Encodes the password into UTF\sphinxhyphen{}8 bytes, generates a random salt,
and creates a secure password hash using bcrypt.
\begin{quote}\begin{description}
\sphinxlineitem{Parameters}
\sphinxAtStartPar
\sphinxstyleliteralstrong{\sphinxupquote{plain\_password}} (\sphinxstyleliteralemphasis{\sphinxupquote{str}}) \textendash{} The original plain\sphinxhyphen{}text password to hash.

\sphinxlineitem{Returns}
\sphinxAtStartPar
The bcrypt password hash decoded as a UTF\sphinxhyphen{}8 string.

\sphinxlineitem{Return type}
\sphinxAtStartPar
str

\end{description}\end{quote}
\subsubsection*{Example}

\sphinxAtStartPar
hashed = hash\_password(“admin123”)
print(hashed)

\end{fulllineitems}

\index{verify\_password() (in module app.security)@\spxentry{verify\_password()}\spxextra{in module app.security}}

\begin{fulllineitems}
\phantomsection\label{\detokenize{modules:app.security.verify_password}}
\pysigstartsignatures
\pysiglinewithargsret
{\sphinxcode{\sphinxupquote{app.security.}}\sphinxbfcode{\sphinxupquote{verify\_password}}}
{\sphinxparam{\DUrole{n}{plain\_password}\DUrole{p}{:}\DUrole{w}{ }\DUrole{n}{str}}\sphinxparamcomma \sphinxparam{\DUrole{n}{password\_hash}\DUrole{p}{:}\DUrole{w}{ }\DUrole{n}{str}}}
{{ $\rightarrow$ bool}}
\pysigstopsignatures
\sphinxAtStartPar
Verify a plain\sphinxhyphen{}text password against an existing bcrypt password hash.

\sphinxAtStartPar
Encodes the provided password and stored hash into UTF\sphinxhyphen{}8 bytes,
then uses bcrypt to check whether the password matches the hash.
\begin{quote}\begin{description}
\sphinxlineitem{Parameters}\begin{itemize}
\item {} 
\sphinxAtStartPar
\sphinxstyleliteralstrong{\sphinxupquote{plain\_password}} (\sphinxstyleliteralemphasis{\sphinxupquote{str}}) \textendash{} The plain\sphinxhyphen{}text password entered by the user.

\item {} 
\sphinxAtStartPar
\sphinxstyleliteralstrong{\sphinxupquote{password\_hash}} (\sphinxstyleliteralemphasis{\sphinxupquote{str}}) \textendash{} The previously stored bcrypt password hash.

\end{itemize}

\sphinxlineitem{Returns}
\sphinxAtStartPar
\begin{description}
\sphinxlineitem{True if the password matches the stored hash;}
\sphinxAtStartPar
otherwise, False.

\end{description}


\sphinxlineitem{Return type}
\sphinxAtStartPar
bool

\end{description}\end{quote}
\subsubsection*{Example}

\sphinxAtStartPar
hashed = hash\_password(“admin123”)
is\_valid = verify\_password(“admin123”, hashed)
print(is\_valid)  \# True

\end{fulllineitems}



\subsection{utils.py}
\label{\detokenize{modules:utils-py}}
\sphinxAtStartPar
The utils.py module contains reusable helper functions.

\sphinxAtStartPar
Responsibilities:
\begin{itemize}
\item {} 
\sphinxAtStartPar
Provide common utility functionality.

\item {} 
\sphinxAtStartPar
Avoid repeating shared logic across application modules.

\end{itemize}

\sphinxAtStartPar
\sphinxstylestrong{Example usage}

\begin{sphinxVerbatim}[commandchars=\\\{\}]
\PYG{k+kn}{from}\PYG{+w}{ }\PYG{n+nn}{app}\PYG{n+nn}{.}\PYG{n+nn}{utils}\PYG{+w}{ }\PYG{k+kn}{import} \PYG{n}{some\PYGZus{}helper}

\PYG{n}{result} \PYG{o}{=} \PYG{n}{some\PYGZus{}helper}\PYG{p}{(}\PYG{l+s+s2}{\PYGZdq{}}\PYG{l+s+s2}{input value}\PYG{l+s+s2}{\PYGZdq{}}\PYG{p}{)}
\end{sphinxVerbatim}

\sphinxAtStartPar
\sphinxstylestrong{API reference}
\index{module@\spxentry{module}!app.utils@\spxentry{app.utils}}\index{app.utils@\spxentry{app.utils}!module@\spxentry{module}}\phantomsection\label{\detokenize{modules:module-app.utils}}
\sphinxAtStartPar
Utility functions for the User Access Control application.

\sphinxAtStartPar
This module provides common helper functions for locating application
resources, loading HTML templates, and retrieving the currently
authenticated user from the session.
\index{current\_user() (in module app.utils)@\spxentry{current\_user()}\spxextra{in module app.utils}}

\begin{fulllineitems}
\phantomsection\label{\detokenize{modules:app.utils.current_user}}
\pysigstartsignatures
\pysiglinewithargsret
{\sphinxcode{\sphinxupquote{app.utils.}}\sphinxbfcode{\sphinxupquote{current\_user}}}
{\sphinxparam{\DUrole{n}{request}\DUrole{p}{:}\DUrole{w}{ }\DUrole{n}{Request}}}
{}
\pysigstopsignatures
\sphinxAtStartPar
Retrieve the currently logged\sphinxhyphen{}in user’s database record.

\sphinxAtStartPar
Reads the user\_id from the request session and uses get\_user\_by\_id()
to retrieve the corresponding user record from the database.

\sphinxAtStartPar
The login function (Member A) stores the user’s ID in the session
after successful authentication.
\begin{quote}\begin{description}
\sphinxlineitem{Parameters}
\sphinxAtStartPar
\sphinxstyleliteralstrong{\sphinxupquote{request}} (\sphinxstyleliteralemphasis{\sphinxupquote{Request}}) \textendash{} The incoming FastAPI request containing
the user’s session information.

\sphinxlineitem{Returns}
\sphinxAtStartPar
The logged\sphinxhyphen{}in user’s database record if a valid
user\_id exists in the session; otherwise, None.

\sphinxlineitem{Return type}
\sphinxAtStartPar
dict | None

\end{description}\end{quote}
\subsubsection*{Example}

\sphinxAtStartPar
user = current\_user(request)
\begin{description}
\sphinxlineitem{if user:}
\sphinxAtStartPar
print(user{[}“username”{]})

\sphinxlineitem{else:}
\sphinxAtStartPar
print(“User is not logged in.”)

\end{description}

\end{fulllineitems}

\index{load\_html() (in module app.utils)@\spxentry{load\_html()}\spxextra{in module app.utils}}

\begin{fulllineitems}
\phantomsection\label{\detokenize{modules:app.utils.load_html}}
\pysigstartsignatures
\pysiglinewithargsret
{\sphinxcode{\sphinxupquote{app.utils.}}\sphinxbfcode{\sphinxupquote{load\_html}}}
{\sphinxparam{\DUrole{n}{file\_name}\DUrole{p}{:}\DUrole{w}{ }\DUrole{n}{str}}}
{{ $\rightarrow$ str}}
\pysigstopsignatures
\sphinxAtStartPar
Read and return the contents of an HTML file from the app/html directory.

\sphinxAtStartPar
Constructs the file path using BASE\_DIR, opens the specified HTML
file with UTF\sphinxhyphen{}8 encoding, and reads its contents.
\begin{quote}\begin{description}
\sphinxlineitem{Parameters}
\sphinxAtStartPar
\sphinxstyleliteralstrong{\sphinxupquote{file\_name}} (\sphinxstyleliteralemphasis{\sphinxupquote{str}}) \textendash{} Name of the HTML file to read.

\sphinxlineitem{Returns}
\sphinxAtStartPar
Contents of the HTML file as a string.

\sphinxlineitem{Return type}
\sphinxAtStartPar
str

\sphinxlineitem{Raises}\begin{itemize}
\item {} 
\sphinxAtStartPar
\sphinxstyleliteralstrong{\sphinxupquote{FileNotFoundError}} \textendash{} If the specified HTML file does not exist.

\item {} 
\sphinxAtStartPar
\sphinxstyleliteralstrong{\sphinxupquote{OSError}} \textendash{} If an error occurs while opening or reading the file.

\end{itemize}

\end{description}\end{quote}
\subsubsection*{Example}

\sphinxAtStartPar
html = load\_html(“login.html”)

\end{fulllineitems}



\subsection{HTML Templates}
\label{\detokenize{modules:html-templates}}
\sphinxAtStartPar
The application uses HTML templates to display its user interface.


\subsubsection{login.html}
\label{\detokenize{modules:login-html}}
\sphinxAtStartPar
Displays the login form.


\subsubsection{home.html}
\label{\detokenize{modules:home-html}}
\sphinxAtStartPar
Displays the home page after successful login.


\subsubsection{denied.html}
\label{\detokenize{modules:denied-html}}
\sphinxAtStartPar
Displays an access denied message.


\subsubsection{page1.html to page4.html}
\label{\detokenize{modules:page1-html-to-page4-html}}
\sphinxAtStartPar
These are protected pages displayed according to user permissions.


\subsection{Static Files}
\label{\detokenize{modules:static-files}}
\sphinxAtStartPar
The static directory contains CSS files used to style the application’s
HTML pages.


\subsection{Conclusion}
\label{\detokenize{modules:conclusion}}
\sphinxAtStartPar
Separating responsibilities into different modules improves readability,
maintainability, and code reuse.

\sphinxstepscope


\section{Testing Documentation}
\label{\detokenize{testing:testing-documentation}}\label{\detokenize{testing::doc}}

\subsection{Introduction}
\label{\detokenize{testing:introduction}}
\sphinxAtStartPar
Testing verifies that the User Access Control application works correctly
and prevents unauthorized access.


\subsection{Authentication Testing}
\label{\detokenize{testing:authentication-testing}}

\subsubsection{Login with Valid Credentials}
\label{\detokenize{testing:login-with-valid-credentials}}\begin{enumerate}
\sphinxsetlistlabels{\arabic}{enumi}{enumii}{}{.}%
\item {} 
\sphinxAtStartPar
Open the login page.

\item {} 
\sphinxAtStartPar
Enter a valid username and password.

\item {} 
\sphinxAtStartPar
Submit the login form.

\item {} 
\sphinxAtStartPar
Verify that the user is redirected to the home page.

\end{enumerate}

\sphinxAtStartPar
Expected Result: The user should log in successfully.


\subsubsection{Login with Invalid Credentials}
\label{\detokenize{testing:login-with-invalid-credentials}}\begin{enumerate}
\sphinxsetlistlabels{\arabic}{enumi}{enumii}{}{.}%
\item {} 
\sphinxAtStartPar
Open the login page.

\item {} 
\sphinxAtStartPar
Enter an incorrect username or password.

\item {} 
\sphinxAtStartPar
Submit the login form.

\end{enumerate}

\sphinxAtStartPar
Expected Result: The application should reject the login attempt.


\subsection{Logout Testing}
\label{\detokenize{testing:logout-testing}}\begin{enumerate}
\sphinxsetlistlabels{\arabic}{enumi}{enumii}{}{.}%
\item {} 
\sphinxAtStartPar
Log in to the application.

\item {} 
\sphinxAtStartPar
Select the logout option.

\item {} 
\sphinxAtStartPar
Try opening a protected page.

\end{enumerate}

\sphinxAtStartPar
Expected Result: The user should be redirected to the login page or denied
access.


\subsection{Authorization Testing}
\label{\detokenize{testing:authorization-testing}}

\subsubsection{Authorized Page Access}
\label{\detokenize{testing:authorized-page-access}}\begin{enumerate}
\sphinxsetlistlabels{\arabic}{enumi}{enumii}{}{.}%
\item {} 
\sphinxAtStartPar
Log in as Priti.

\item {} 
\sphinxAtStartPar
Open Page 1 or Page 2.

\end{enumerate}

\sphinxAtStartPar
Expected Result: Access should be granted.


\subsubsection{Unauthorized Page Access}
\label{\detokenize{testing:unauthorized-page-access}}\begin{enumerate}
\sphinxsetlistlabels{\arabic}{enumi}{enumii}{}{.}%
\item {} 
\sphinxAtStartPar
Log in as Priti.

\item {} 
\sphinxAtStartPar
Try opening Page 3.

\end{enumerate}

\sphinxAtStartPar
Expected Result: Access should be denied.


\subsubsection{Administrator Access}
\label{\detokenize{testing:administrator-access}}\begin{enumerate}
\sphinxsetlistlabels{\arabic}{enumi}{enumii}{}{.}%
\item {} 
\sphinxAtStartPar
Log in as Admin.

\item {} 
\sphinxAtStartPar
Open Page 1, Page 2, Page 3, and Page 4.

\end{enumerate}

\sphinxAtStartPar
Expected Result: Access should be granted to all four pages.


\subsection{Session Testing}
\label{\detokenize{testing:session-testing}}
\sphinxAtStartPar
Verify that:
\begin{itemize}
\item {} 
\sphinxAtStartPar
A logged\sphinxhyphen{}in user’s session is maintained while navigating pages.

\item {} 
\sphinxAtStartPar
The session expires according to the configured expiration time.

\item {} 
\sphinxAtStartPar
Logout clears the session.

\end{itemize}


\subsection{Database Testing}
\label{\detokenize{testing:database-testing}}
\sphinxAtStartPar
Verify that:
\begin{itemize}
\item {} 
\sphinxAtStartPar
User records can be retrieved from MySQL.

\item {} 
\sphinxAtStartPar
Password hashes are stored correctly.

\item {} 
\sphinxAtStartPar
User page permissions are retrieved correctly.

\end{itemize}


\subsection{Conclusion}
\label{\detokenize{testing:conclusion}}
\sphinxAtStartPar
Testing helps identify errors and ensures that authentication, authorization,
session management, and database operations work as expected.

\sphinxstepscope


\section{Deployment Guide}
\label{\detokenize{deployment:deployment-guide}}\label{\detokenize{deployment::doc}}

\subsection{Introduction}
\label{\detokenize{deployment:introduction}}
\sphinxAtStartPar
This guide explains the basic steps required to run the User Access Control
application in a local or server environment.


\subsection{Prerequisites}
\label{\detokenize{deployment:prerequisites}}
\sphinxAtStartPar
Make sure the following software is installed:
\begin{itemize}
\item {} 
\sphinxAtStartPar
Python 3.13

\item {} 
\sphinxAtStartPar
MySQL Server

\item {} 
\sphinxAtStartPar
Required Python packages

\item {} 
\sphinxAtStartPar
A compatible operating system

\end{itemize}


\subsection{Project Configuration}
\label{\detokenize{deployment:project-configuration}}\begin{enumerate}
\sphinxsetlistlabels{\arabic}{enumi}{enumii}{}{.}%
\item {} 
\sphinxAtStartPar
Download or clone the project repository.

\item {} 
\sphinxAtStartPar
Open the project directory.

\item {} 
\sphinxAtStartPar
Configure the database connection.

\item {} 
\sphinxAtStartPar
Set the required environment variables.

\item {} 
\sphinxAtStartPar
Install the Python dependencies.

\end{enumerate}


\subsection{Database Setup}
\label{\detokenize{deployment:database-setup}}\begin{enumerate}
\sphinxsetlistlabels{\arabic}{enumi}{enumii}{}{.}%
\item {} 
\sphinxAtStartPar
Start the MySQL server.

\item {} 
\sphinxAtStartPar
Create the required database.

\item {} 
\sphinxAtStartPar
Create the required database tables.

\item {} 
\sphinxAtStartPar
Configure database credentials.

\item {} 
\sphinxAtStartPar
Initialize the application data if required.

\end{enumerate}


\subsection{Run the Application Locally}
\label{\detokenize{deployment:run-the-application-locally}}
\sphinxAtStartPar
Open the terminal in the project directory and execute:

\begin{sphinxVerbatim}[commandchars=\\\{\}]
python\PYG{+w}{ }\PYGZhy{}m\PYG{+w}{ }uvicorn\PYG{+w}{ }app.main:app\PYG{+w}{ }\PYGZhy{}\PYGZhy{}reload
\end{sphinxVerbatim}

\sphinxAtStartPar
The application will be available at:

\begin{sphinxVerbatim}[commandchars=\\\{\}]
http://127.0.0.1:8000
\end{sphinxVerbatim}


\subsection{Production Deployment}
\label{\detokenize{deployment:production-deployment}}
\sphinxAtStartPar
For production deployment, consider the following practices:
\begin{itemize}
\item {} 
\sphinxAtStartPar
Use HTTPS.

\item {} 
\sphinxAtStartPar
Store secrets in environment variables.

\item {} 
\sphinxAtStartPar
Disable development reload mode.

\item {} 
\sphinxAtStartPar
Use a production ASGI server configuration.

\item {} 
\sphinxAtStartPar
Configure secure session cookies.

\item {} 
\sphinxAtStartPar
Restrict database access.

\item {} 
\sphinxAtStartPar
Configure application logging.

\item {} 
\sphinxAtStartPar
Back up the database regularly.

\item {} 
\sphinxAtStartPar
Monitor application errors and availability.

\end{itemize}


\subsection{Security Recommendations}
\label{\detokenize{deployment:security-recommendations}}\begin{itemize}
\item {} 
\sphinxAtStartPar
Never commit passwords or secret keys to Git.

\item {} 
\sphinxAtStartPar
Use strong session secret keys.

\item {} 
\sphinxAtStartPar
Keep dependencies updated.

\item {} 
\sphinxAtStartPar
Use secure database credentials.

\item {} 
\sphinxAtStartPar
Apply appropriate access permissions.

\end{itemize}


\subsection{Conclusion}
\label{\detokenize{deployment:conclusion}}
\sphinxAtStartPar
Following these steps helps prepare the User Access Control application
for deployment in a suitable environment.


\renewcommand{\indexname}{Python Module Index}
\begin{sphinxtheindex}
\let\bigletter\sphinxstyleindexlettergroup
\bigletter{a}
\item\relax\sphinxstyleindexentry{app.access}\sphinxstyleindexpageref{modules:\detokenize{module-app.access}}
\item\relax\sphinxstyleindexentry{app.auth}\sphinxstyleindexpageref{modules:\detokenize{module-app.auth}}
\item\relax\sphinxstyleindexentry{app.db}\sphinxstyleindexpageref{modules:\detokenize{module-app.db}}
\item\relax\sphinxstyleindexentry{app.main}\sphinxstyleindexpageref{modules:\detokenize{module-app.main}}
\item\relax\sphinxstyleindexentry{app.security}\sphinxstyleindexpageref{modules:\detokenize{module-app.security}}
\item\relax\sphinxstyleindexentry{app.utils}\sphinxstyleindexpageref{modules:\detokenize{module-app.utils}}
\end{sphinxtheindex}

\renewcommand{\indexname}{Index}
\printindex
\end{document}